\section{Validation Framework}

City1 v2 uses a multi-layer validation framework because no single benchmark can close all methodological risks in this problem setting. Different evidence layers answer different questions about the credibility and usefulness of the final proxy population surface, and the manuscript treats them as complementary rather than interchangeable. Supplementary Table~S1 makes this evidentiary logic explicit by mapping each layer to the question it answers, the artifact it uses, and the claims it does not support.

Leave-one-city-out validation addresses transferability to unseen cities and is therefore central to the baseline claim of generalization under data scarcity. In this protocol, one city is held out as an unseen evaluation target while the remaining cities are used for training; this process is repeated across the available city set and summarized across the hold-out runs. The purpose of this layer is not to maximize a single pooled score, but to ask whether the baseline remains usable once it is required to transfer beyond the cities seen during training.

Spatial block cross-validation addresses a different concern: whether performance is sustained once local spatial leakage within a city is restricted. Rather than testing transfer to a new city, this layer is used as a stricter within-city robustness check. Taken together, leave-one-city-out and spatial block cross-validation help distinguish inter-city transferability from within-city robustness instead of collapsing both issues into a single performance summary.

The framework is extended by a grid-size benchmark, which compares 250\,m, 500\,m, and 1000\,m configurations and is used to justify the frozen production resolution rather than as a cosmetic sensitivity analysis. Its function is to support the production choice of 500\,m as the frozen operating grid, not to present an open-ended sensitivity theater in which every resolution is treated as equally viable. OSM completeness serves as an input-reliability layer rather than as a direct measure of predictive accuracy. Its purpose is to contextualize variation in feature support across cities and to inform how strongly later evidence layers can be interpreted.

District benchmark comparison provides partial internal administrative validation by aggregating cell-level predictions to districts and comparing them against district-level references where available. This layer is useful, but it is not presented as solved ground-truth validation. It tests whether the proxy surface remains broadly coherent after aggregation to a coarser administrative support, while also revealing where this coherence is heterogeneous across cities.

External benchmark comparison serves a different role: it evaluates structural agreement with independent gridded population products and is therefore interpreted as external structural comparison rather than as comparison against cell-level truth. In the frozen package, this layer is summarized through Pearson correlation, Spearman correlation, top-decile overlap, and hotspot IoU. These statistics are intended to compare broad ranking agreement and dense concentration structure, not to suggest that either external product provides the true population count per grid cell.

Additional layers refine the interpretation of the model from other directions. Ablation isolates the contribution of major feature families and clarifies whether performance depends primarily on built-form structure or on the broader feature stack. Qualitative validation examines whether the resulting surfaces are spatially plausible through curated hotspot and coldspot cases. In the frozen manuscript package, this qualitative layer is locked to Almaty and Astana because they provide two contrasting reading contexts: a stronger completeness case and a more cautious moderate-completeness case. The goal is not cell-level validation, but plausibility review through deliberately selected hotspot and coldspot examples.

In combination, the validation framework is intended to support a bounded baseline claim grounded in multiple evidence layers rather than in any single metric or benchmark. The framework is therefore less a hierarchy with one decisive test at the top than a structured set of checks that examine transferability, within-city robustness, internal administrative coherence, external structural agreement, feature dependence, and map plausibility from different methodological angles.
